\chapter{Glosario de términos}
\label{chap:glosario-terminos}

\begin{description}

  \item[Aceleración (\emph{speedup})] Cociente entre el tiempo de ejecución
        con un hilo y el tiempo con $t$ hilos dentro del mismo modo y clase.
        Un valor superior a~1 indica ganancia con el paralelismo; un valor
        inferior a~1 indica degradación.

  \item[Barrera (\emph{barrier})] Punto de sincronización que obliga a todos
        los hilos de una región paralela a esperar hasta que el último de
        ellos llegue al mismo punto antes de continuar.

  \item[Benchmark] Programa de referencia diseñado para medir el rendimiento
        de un sistema de forma reproducible y comparable entre plataformas.

  \item[Cython] Compilador que acepta un superconjunto del lenguaje Python
        con anotaciones de tipo opcionales y genera código C que se compila
        como módulo de extensión para CPython. Las anotaciones de tipo
        permiten sustituir objetos Python por tipos C nativos en los bucles
        críticos.

  \item[\emph{Free-threaded Python}] Variante experimental de CPython,
        disponible desde la versión~3.13 mediante la opción
        \texttt{--disable-gil}, que elimina el GIL y permite la ejecución
        simultánea de múltiples hilos sobre el bytecode del intérprete.

  \item[Kernel numérico] Función que concentra la mayor parte del cómputo
        útil de un benchmark o aplicación paralela, separada del código de
        orquestación (inicialización, E/S, verificación).

  \item[Memoryview] Objeto Python que expone el buffer interno de un array
        NumPy sin copiar los datos. En Cython, las memoryviews tipadas
        permiten acceder directamente a los elementos del array como punteros
        C, eliminando la indirección de los objetos Python en los bucles
        internos.

  \item[Modo de ejecución] En OMP4Py, configuración que determina cómo se
        procesan las directivas OpenMP: modo~0 (puro, sin transformación),
        modo~1 (híbrido, con transformación AST pero sin compilación),
        modo~2 (compilado con Cython) y modo~3 (compilado con tipado
        Cython).

  \item[OMP4Py] Implementación en Python puro del modelo de programación
        OpenMP. Traduce directivas OpenMP embebidas como cadenas de texto en
        código Python paralelo mediante transformación del AST, con soporte
        opcional de compilación Cython.

  \item[Overhead de sincronización] Tiempo invertido en operaciones de
        coordinación entre hilos (barreras, secciones críticas, reducciones)
        que no contribuye directamente al cómputo útil.

  \item[Reducción] Operación que combina los valores parciales calculados por
        cada hilo en un único resultado escalar (suma, máximo, mínimo, etc.)
        al final de una región paralela.

  \item[Región paralela] Bloque de código delimitado por una directiva
        \texttt{parallel} de OpenMP que se ejecuta simultáneamente en
        múltiples hilos.

  \item[Slurm] Sistema de gestión y planificación de trabajos utilizado en
        los supercomputadores del CESGA para asignar recursos de cómputo y
        organizar la ejecución de trabajos en cola.

  \item[Tipado Cython] Proceso de anotar variables en código Cython con tipos
        C explícitos (\texttt{cython.int}, \texttt{cython.double}, etc.) para
        que el compilador genere operaciones nativas en lugar de pasar por el
        mecanismo de despacho de objetos Python.

  \item[Verificación NAS] Comprobación numérica integrada en cada NAS
        Parallel Benchmark que calcula un checksum de los resultados y lo
        compara con un valor de referencia para la clase ejecutada.
        Una verificación \texttt{SUCCESSFUL} garantiza la corrección del
        cómputo.

\end{description}
